\section{Convex star-shaped projections and exact formulas for the weak exponent}\label{sec:gauges}
This section proves that the projection of every cycle to the $(U,p)$ plane is strictly star-shaped about the origin and strictly convex. The two positive functions that express these properties satisfy linear equations driven by the convexity gauges $K$ and $W$, and averaging those equations gives exact formulas for the weak exponent $\lambda_c$. The remaining subsections record the consequences: a reduction of the adjoint ratios to one scalar Riccati equation, a polar form, a comparison of nested cycles, sign-definite integrals, exact moment identities, a pairing with the strong-stable Floquet solution, and injectivity of the Gauss map. Besides the identities of Section~\ref{sec:adjoint}, the proofs use Theorems~\ref{thm:support} and~\ref{thm:geometry}, Corollary~\ref{cor:cones} and Lemma~\ref{lem:curvature}; Proposition~\ref{prop:strongidentity} also uses Lemma~\ref{lem:strongstable} and the compound system \eqref{eq:compound}. No proof uses the stability of a cycle, the number of cycles, or any property of $(\mu,k)$ other than the existence of a cycle. Floating-point evidence for the formulas is reported in Section~\ref{sec:evidence}; it plays no part in the proofs.

\paragraph{Local notation.}
A cycle with least period $T$ is fixed unless two cycles are named. By Theorem~\ref{thm:geometry}, $p$ has exactly two zeros in a period, both simple: the \emph{trough}, at a time $t_a$, where $U=U_{\min}<0$, and the \emph{crest}, at the time $t_b\in(t_a,t_a+T)$, where $U=U_{\max}>0$. The \emph{rising arc} is $(t_a,t_b)$, on which $p>0$; the \emph{falling arc} is $(t_b,t_a+T)$, on which $p<0$. We write $q=(U,p)$ for the projected point, $v\wedge w=v_1w_2-v_2w_1$ for planar vectors, and
\[
L_2=D^2+(\mu-\eta)D+\zeta
\]
for the second-order factor in \eqref{eq:factorisation}. Besides the global symbols $\Sg$, $\Xg$, $\kappa=K/\Sg$ and $\etaz$, this section uses: $\zeta_0=\Xg/\Sg$ and $\eta_1$ (Theorem~\ref{thm:gauges}); $\Sg_c=p^2-(U+c)p'$ for a shift $c\ge0$; polar coordinates $(U,p)=r(\cos\theta,\sin\theta)$ with $s=(y-k)/r$ and $\sigma=\Sg/r^2$; for two nested cycles $R_i(\theta)$, $s_i(\theta)$, $\sigma_i(\theta)$, $i=1,2$; the chain ratios $r_1=x/y$, $r_2=y/\ell$; the polynomials $f_\mu(s)=s(s+1)(s+1-\mu)$ of \eqref{eq:gainprimitive}, $\chi(s)=s(s-1)(s-\mu)$ and $G_\mu(s)=\tfrac32s^2+(2-\mu)s$; the operator $\mathcal L[w]=w''+(3w+2-\mu)w'+f_\mu(w)$; the variable $\omega=\zeta+\eta+\eta^2/2$ (in Proposition~\ref{prop:gaussmap} only); and a hat for the periodic part of a strong-stable Floquet component, $\hat a=e^{\gamma t}a$. The letters $r_1,r_2$ denote chain ratios and are unrelated to the polar radius $r$.

\subsection{Two positive gauges}
\begin{lemma}[Wronskian gauge identity]\label{gau:wronskian}
Let $f,g$ be $C^2$ functions of $t$ along a cycle and put $V_{f,g}=fg'-f'g$. Then
\begin{equation}\label{gau:wronskianeq}
V_{f,g}'=(\eta-\mu)V_{f,g}+f\,L_2[g]-g\,L_2[f].
\end{equation}
Moreover
\begin{equation}\label{gau:Ltwovalues}
L_2[1]=\zeta,\qquad L_2[U]=K,\qquad L_2[p]=0,\qquad L_2[-p']=W .
\end{equation}
\end{lemma}
\begin{proof}
$V_{f,g}'=fg''-f''g$. Write $g''=L_2[g]-(\mu-\eta)g'-\zeta g$ and the same for $f$; the terms $\zeta fg$ cancel and
\[
fg''-f''g=f\,L_2[g]-g\,L_2[f]+(\eta-\mu)(fg'-f'g).
\]
The first value in \eqref{gau:Ltwovalues} is immediate. Since $U'=p$ and $p'=k-\mu p-y$,
\[
L_2[U]=p'+(\mu-\eta)p+\zeta U=k-y-\eta p+\zeta U=K .
\]
The identity $L_2[p]=0$ is \eqref{eq:tangent}. Differentiating it gives $L_2[p']-\eta'p'+\zeta'p=0$, that is, $L_2[-p']=\zeta'p-\eta'p'=W$ by \eqref{eq:curvature}.
\end{proof}

\begin{theorem}[Star gauge, convexity gauge and exact exponent]\label{thm:gauges}
On every cycle of \eqref{eq:logsystem} the following hold.
\begin{enumerate}[label=(\roman*)]
\item \emph{Positivity.} The functions $\Sg=p^2-Up'$ and $\Xg=p'^2-pp''$ are positive at every phase, and
\begin{equation}\label{eq:gaugeodes}
\Sg'=(\eta-\mu)\Sg+pK,\qquad \Xg'=(\eta-\mu)\Xg+pW .
\end{equation}
Both are functions of the state: with $p'=k-\mu p-y$, one has $\Sg=p^2-Up'$ and $\Xg=p'(k-y)+p(x-y)$.
\item \emph{Star shape.} The projected cycle does not pass through the origin, its polar angle satisfies
\[
\frac{\mathrm d}{\mathrm dt}\arg(U+ip)=-\frac{\Sg}{U^2+p^2}<0
\]
and decreases by exactly $2\pi$ in one period. Every ray from the origin meets the projected cycle in exactly one point, and transversally.
\item \emph{Strict convexity.} The signed curvature of the projected cycle is $-\Xg/(p^2+p'^2)^{3/2}<0$, and its tangent angle $\arg(p+ip')$ decreases by exactly $2\pi$ in one period. Every tangent line is a strict supporting line: for every $t_1$ and every point $q\ne q(t_1)$ of the closed region bounded by the projected cycle, in particular for $q=q(t)$ with $t\not\equiv t_1\pmod T$ and for $q=(0,0)$,
\begin{equation}\label{gau:supportline}
\bigl(p(t_1),p'(t_1)\bigr)\wedge\bigl(q-q(t_1)\bigr)<0 .
\end{equation}
That region is the convex hull of the projected cycle.
\item \emph{Exact exponent.} With $\oint f\,\mathrm dU=\int_0^Tf\,p\dd$,
\begin{equation}\label{eq:exactexponent}
\lambda_c=-\avg{\frac{pK}{\Sg}}=-\avg{\frac{pW}{\Xg}},\qquad
\lambda_cT=-\oint\frac{K}{\Sg}\,\mathrm dU=-\oint\frac{W}{\Xg}\,\mathrm dU .
\end{equation}
More generally, for every $\vartheta\in[0,1]$,
\begin{equation}\label{gau:mixedexponent}
\lambda_c=-\avg{\frac{p\,[(1-\vartheta)K+\vartheta W]}{(1-\vartheta)\Sg+\vartheta\,\Xg}} .
\end{equation}
\item \emph{Shifted centres.} For every $c$ with $0\le c<|U_{\min}|$ the function $\Sg_c=p^2-(U+c)p'$ is positive at every phase,
\begin{equation}\label{gau:shifted}
\Sg_c'=(\eta-\mu)\Sg_c+p\,(K+c\,\zeta),\qquad
\lambda_c=-\avg{\frac{p\,(K+c\,\zeta)}{\Sg_c}},
\end{equation}
and the conclusions of \textup{(ii)} hold for rays from $(-c,0)$.
\item \emph{Explicit comparison functions.} The state functions
\begin{equation}\label{gau:etazero}
\etaz=\frac{U(x-y)-p(y-k)}{\Sg},\qquad
\eta_1=\frac{x\Sg+k\psiU(U)\,p'}{\Xg}-1
\end{equation}
satisfy $(\log\Sg)'=\etaz-\mu$ and $(\log\Xg)'=\eta_1-\mu$, hence $\avg{\etaz}=\avg{\eta_1}=\mu$, and
\begin{equation}\label{gau:etadifference}
\eta-\etaz=-\frac{pK}{\Sg},\qquad \eta-\eta_1=-\frac{pW}{\Xg}.
\end{equation}
Consequently $\eta<\min(\etaz,\eta_1)$ on the rising arc, $\eta>\max(\etaz,\eta_1)$ on the falling arc, $\eta=\etaz=\eta_1$ at the crest and at the trough, and $\lambda_c=\avg{\eta-\etaz}=\avg{\eta-\eta_1}$.
\end{enumerate}
\end{theorem}
\begin{proof}
\emph{The gauge equations.} In Lemma~\ref{gau:wronskian} take $f=p$, so that $L_2[f]=0$. For $g=U$ one has $V_{p,U}=pU'-p'U=\Sg$ and $L_2[U]=K$. For $g=-p'$ one has $V_{p,-p'}=p'^2-pp''=\Xg$ and $L_2[-p']=W$. For $g=U+c$ one has $V_{p,U+c}=\Sg-cp'=\Sg_c$ and $L_2[U+c]=K+c\zeta$. Equation \eqref{gau:wronskianeq} gives \eqref{eq:gaugeodes} and the first equation in \eqref{gau:shifted}. Since $p''=-\mu p'-(x-y)$ and $p'+\mu p=k-y$,
\[
\Xg=p'^2+\mu pp'+p(x-y)=p'(k-y)+p(x-y).
\]

\emph{Values at the zeros of $p$.} At a zero of $p$, \eqref{eq:crestidentities} gives $x-k=(1+\eta)(y-k)$ and $p'=k-y$. If $U=0$ there, then $x=k$, hence $y=k$ and $p'=0$, and the point would be the equilibrium. Thus $U\ne0$, and $y-k$ has the sign of $x-k$, which is the sign of $U$. At $t_a$ and $t_b$ therefore
\[
\Sg=-Up'=U(y-k)>0,\qquad \Xg=p'^2=(y-k)^2>0,\qquad \Sg_c=(U+c)(y-k).
\]
At the crest $U+c>0$ and $y>k$. At the trough $U+c=U_{\min}+c<0$ and $y<k$. Hence $\Sg_c>0$ at both zeros when $0\le c<|U_{\min}|$.

\emph{Positivity.} Let $(S,N)$ be one of the pairs $(\Sg,K)$, $(\Xg,W)$, $(\Sg_c,K+c\zeta)$. Then $S'=(\eta-\mu)S+pN$, $S(t_a)>0$, and $N>0$ on the whole cycle by \eqref{eq:positiveK}, Lemma~\ref{lem:curvature} and $\zeta>0$. Put $\mathcal E(t)=\exp\int_{t_a}^t(\eta-\mu)\dd$. Then
\[
\Bigl(\frac{S}{\mathcal E}\Bigr)'=\frac{pN}{\mathcal E}.
\]
On $[t_a,t_b]$ we have $p\ge0$, so $S/\mathcal E$ is nondecreasing and $S(t)\ge\mathcal E(t)S(t_a)>0$. On $[t_b,t_a+T]$ we have $p\le0$, so $S/\mathcal E$ is nonincreasing and
\[
\frac{S(t)}{\mathcal E(t)}\ge\frac{S(t_a+T)}{\mathcal E(t_a+T)}=\frac{S(t_a)}{\mathcal E(t_a+T)}>0 .
\]
The two arcs exhaust the period. This proves (i) and the positivity in (v).

\emph{Star shape.} Since $\Sg>0$, the point $(U,p)$ is never the origin, and a continuous branch $\theta(t)$ of $\arg(U+ip)$ has
\[
\theta'=\frac{Up'-pU'}{U^2+p^2}=-\frac{\Sg}{U^2+p^2}<0 .
\]
By periodicity $\theta(t+T)-\theta(t)=-2\pi n$ with an integer $n\ge1$. A strictly decreasing continuous function maps a half-open period bijectively onto a half-open interval of length $2\pi n$, which contains exactly $2n$ multiples of $\pi$. The angle is a multiple of $\pi$ exactly when $p=0$, which happens twice in a period. Hence $n=1$, and $\theta$ is a bijection from a half-open period onto a half-open interval of length $2\pi$. This is (ii); transversality is $\theta'\ne0$. The same argument with $\arg(U+c+ip)$, whose derivative is $-\Sg_c/((U+c)^2+p^2)$, gives the last assertion of (v).

\emph{Strict convexity.} The projected velocity is $(p,p')$ and the acceleration is $(p',p'')$, so the signed curvature is $(pp''-p'^2)/(p^2+p'^2)^{3/2}$. The projection is regular by Theorem~\ref{thm:geometry}, so a continuous branch $\phi(t)$ of $\arg(p+ip')$ exists, and $\phi'=-\Xg/(p^2+p'^2)<0$. The angle $\phi$ lies in $\pi/2+\pi\mathbb Z$ exactly when $p=0$, so the counting argument just given shows that $\phi$ decreases by exactly $2\pi$ in a period. Consequently, for every unit vector $\nu$ the velocity $(p,p')$ is orthogonal to $\nu$ at exactly two phases in a period.

Fix $t_1$, put $n_1=(-p'(t_1),p(t_1))$ and $h(t)=n_1\cdot(q(t)-q(t_1))$, so that $h(t)$ is the left side of \eqref{gau:supportline} for $q=q(t)$. Then $h'(t)=n_1\cdot(p(t),p'(t))$ vanishes at exactly two phases in a period, one of which is $t_1$. A nonconstant periodic $C^1$ function with exactly two critical phases attains its maximum at one of them and its minimum at the other, and at no other phase. Since $h(t_1)=0$, $h'(t_1)=0$ and $h''(t_1)=p(t_1)p''(t_1)-p'(t_1)^2=-\Xg(t_1)<0$, the phase $t_1$ is not the minimum. It is therefore the maximum, and $h(t)<0$ for $t\not\equiv t_1\pmod T$.

The bounded complementary component of a Jordan curve lies in the convex hull of the curve, because a point outside the hull is joined to infinity by a ray that does not meet the hull. Thus every point $q$ of the closed region bounded by the projected cycle is a convex combination of points $q(t_j)$ of the curve. The affine function $q\mapsto n_1\cdot(q-q(t_1))$ is negative at each $q(t_j)\ne q(t_1)$ and zero at $q(t_1)$, so it is negative at $q$ unless $q=q(t_1)$. This proves \eqref{gau:supportline}; the origin lies in the region by Theorem~\ref{thm:geometry}. Each point of the curve lies on a supporting line of the hull, hence on the boundary of the hull. The boundary of a compact convex set with nonempty interior is a Jordan curve, and a Jordan curve contained in another Jordan curve coincides with it. The projected cycle is therefore the boundary of its convex hull. The interior of the hull is connected, does not meet the curve and contains the bounded complementary component, so it coincides with that component; this proves the last sentence of (iii).

\emph{Exact exponent.} Let $(S,N)$ be as above, or $S=(1-\vartheta)\Sg+\vartheta\,\Xg$ and $N=(1-\vartheta)K+\vartheta W$. Then $S$ is positive and $T$-periodic with $(\log S)'=\eta-\mu+pN/S$. The mean of $(\log S)'$ over a period is zero, and $\lambda_c=\avg\eta-\mu$ by \eqref{eq:exponent}. This gives \eqref{eq:exactexponent}, \eqref{gau:mixedexponent} and the second formula in \eqref{gau:shifted}. The line-integral forms use $\mathrm dU=p\,\mathrm dt$.

\emph{Comparison functions.} From $\Sg'=pp'-Up''$ and $p''=-\mu p'-(x-y)$,
\[
\Sg'=pp'+\mu Up'+U(x-y)=-\mu\,(p^2-Up')+U(x-y)+p\,(p'+\mu p)=-\mu\Sg+U(x-y)-p(y-k),
\]
because $p'+\mu p=k-y$. Hence $(\log\Sg)'=\etaz-\mu$. Next, $\Xg'=p'p''-pp'''$ and $p'''=-\mu p''-xp+(x-y)$, so
\[
\Xg'=(p'+\mu p)\,p''+xp^2-p(x-y).
\]
On the other hand $x\Sg-A\Xg+k\psiU(U)p'=xp^2+A(pp''-p'^2)-(x-k)p'$, since $k\psiU(U)=xU-x+k$. The difference of the two expressions is
\[
p'\bigl[p''+Ap'+x-k\bigr]-p\bigl[p''+(x-y)\bigr]=p'\cdot(-\mu p)-p\cdot(-\mu p')=0,
\]
where $p''+Ap'+x-k=p'+y-k=-\mu p$ and $p''+(x-y)=-\mu p'$. Thus
\begin{equation}\label{gau:Xistate}
\Xg'=x\Sg-A\Xg+k\psiU(U)\,p',
\end{equation}
which is $(\log\Xg)'=\eta_1-\mu$. Both logarithms are periodic, so $\avg{\etaz}=\avg{\eta_1}=\mu$. Subtracting these state identities from \eqref{eq:gaugeodes} divided by $\Sg$ and by $\Xg$ gives \eqref{gau:etadifference}. The sign statements follow from $K,W,\Sg,\Xg>0$ and the sign of $p$.
\end{proof}

The functions $\etaz$ and $\eta_1$ need not be positive. In floating-point computations $\eta_1$ takes negative values on part of every cycle examined, and $\etaz$ on part of some cycles with small $\mu$ (Section~\ref{sec:evidence}).

\begin{corollary}[Half-cycle bounds]\label{gau:halfcycle}
Let $I_+=\int_{t_a}^{t_b}(\eta-\mu)\dd$ and $I_-=\int_{t_b}^{t_a+T}(\eta-\mu)\dd$, so that $I_++I_-=\lambda_cT$, and let $y_a<k<y_b$ be the values of $y$ at the trough and at the crest. Then
\[
I_+<\log\frac{S(t_b)}{S(t_a)},\qquad I_->-\log\frac{S(t_b)}{S(t_a)}\qquad\text{for }S=\Sg\text{ and for }S=\Xg,
\]
where
\[
\frac{\Sg(t_b)}{\Sg(t_a)}=\frac{U_{\max}(y_b-k)}{U_{\min}(y_a-k)},\qquad
\frac{\Xg(t_b)}{\Xg(t_a)}=\frac{(y_b-k)^2}{(k-y_a)^2}.
\]
\end{corollary}
\begin{proof}
With $\mathcal E$ and $(S,N)$ as in the proof of Theorem~\ref{thm:gauges}, $(S/\mathcal E)'=pN/\mathcal E$ is positive on the rising arc and negative on the falling arc. Hence $S(t_b)/\mathcal E(t_b)>S(t_a)$, which is $e^{I_+}<S(t_b)/S(t_a)$, and $S(t_a+T)/\mathcal E(t_a+T)<S(t_b)/\mathcal E(t_b)$, which is $e^{I_-}>S(t_a)/S(t_b)$. The values of $S$ at the zeros of $p$ were computed in that proof.
\end{proof}
The upper bound for the rising arc and the lower bound for the falling arc are opposite numbers; their sum contains no information about $\lambda_cT$.

\begin{remark}[Further readings of the gauges]\label{gau:remarks}
(a) Since $(p/U)'=-\Sg/U^2$ and $(p'/p)'=-\Xg/p^2$, the function $|U|=|\log X|$ is strictly log-concave on each of the two arcs on which $U\ne0$, and $|p|=\lambda|1-Y|$ is strictly log-concave on the rising arc and on the falling arc.

(b) By \eqref{gau:wronskianeq}, whenever $f$ and $g$ are $T$-periodic and $V_{f,g}>0$ on a cycle,
\[
\lambda_c=\avg{\frac{g\,L_2[f]-f\,L_2[g]}{V_{f,g}}} .
\]
Theorem~\ref{thm:gauges} is the case $f=p$, in which $L_2[f]=0$ and the remaining factor $L_2[g]\in\{K,W,K+c\zeta\}$ is positive.

(c) On the cycle, $\Sg$ is the numerator of the angular velocity and $\Xg$ the numerator of the curvature of the planar projection. Equations \eqref{eq:gaugeodes} state that $1/\Sg$ and $1/\Xg$ are Dulac-type weights whose defects are $pK$ and $pW$. The defects change sign with $p$, so the sign hypothesis of the extended Bendixson--Dulac criterion built on the curvature numerator \cite[Theorem~2.3]{GG21} does not hold; see Section~\ref{sec:negative}.
\end{remark}

\subsection{Reduction to one scalar unknown}
\begin{proposition}[Scalar Riccati reduction]\label{prop:kappa}
On every cycle let $\kappa=K/\Sg>0$ and $\zeta_0=\Xg/\Sg>0$.
\begin{enumerate}[label=(\roman*)]
\item \emph{Affine parametrisation.} The adjoint ratios and both convexity gauges are affine in $\kappa$ with coefficients that are functions of the state:
\begin{equation}\label{gau:affine}
\eta=\etaz-p\,\kappa,\qquad \zeta=\zeta_0-p'\kappa,\qquad K=\kappa\,\Sg,\qquad W=W_0+\kappa\,\Xg,
\end{equation}
where $W_0=p\,[x-(1+\etaz)\zeta_0]-p'\,[\zeta_0+(\mu-1)\etaz-\etaz^2]$. Moreover $\zeta_0p-\etaz p'=x-y$ and $\zeta_0U-\etaz p-y+k=0$.
\item \emph{Riccati equation.} The function $\kappa$ is a positive $T$-periodic solution of
\begin{equation}\label{eq:kappariccati}
\kappa'=\frac{k\psiU(U)}{\Sg}-(1+2\etaz-\mu)\,\kappa+p\,\kappa^2 ,
\end{equation}
whose coefficients are functions of the state; equivalently $\kappa'=k\psiU(U)/\Sg-(1+\eta+\etaz-\mu)\kappa$.
\item \emph{Compound form.} The vector $\omega_0=(\Sg,\ \Sg\etaz,\ \Xg)^{\mathsf T}$ satisfies, with $C$ as in \eqref{eq:compound},
\begin{equation}\label{gau:compoundsource}
\omega_0'=C\omega_0+k\psiU(U)\,(0,p,p')^{\mathsf T}.
\end{equation}
Equivalently, $(\etaz,\zeta_0)$ solves the projective system \eqref{eq:projective} up to explicit defects:
\[
\etaz'=\zeta_0+(\mu-1)\etaz-\etaz^2+\frac{k\psiU(U)\,p}{\Sg},\qquad
\zeta_0'=x-(1+\etaz)\zeta_0+\frac{k\psiU(U)\,p'}{\Sg}.
\]
\item \emph{Branch form of the exponent.} For $U\in(U_{\min},U_{\max})$ let $\kappa_{\rm up}(U)$ and $\kappa_{\rm down}(U)$ be the values of $\kappa$ at the point of the rising arc and at the point of the falling arc with abscissa $U$. Then
\begin{equation}\label{gau:branch}
\lambda_cT=-\int_{U_{\min}}^{U_{\max}}\bigl[\kappa_{\rm up}(U)-\kappa_{\rm down}(U)\bigr]\,\mathrm dU .
\end{equation}
\item \emph{Certificates by state functions.} Let $\Lambda$ be a $C^1$ function of the state $(U,p,y)$ on a neighbourhood of the cycle and put $\varepsilon_\Lambda=\Dt\Lambda+\etaz-\mu$, a function of the state. Then along the cycle
\begin{equation}\label{gau:certificate}
\Dt\Lambda+\eta-\mu=\varepsilon_\Lambda-p\,\kappa,\qquad \avg{\Dt\Lambda+\eta-\mu}=\lambda_c .
\end{equation}
Hence $\Dt\Lambda+\eta-\mu\le0$ on the cycle if and only if
\begin{gather*}
\kappa\ge\frac{\varepsilon_\Lambda}{p}\ \text{ on the rising arc},\qquad
\kappa\le\frac{\varepsilon_\Lambda}{p}\ \text{ on the falling arc},\\
\varepsilon_\Lambda\le0\ \text{ at the crest and at the trough}.
\end{gather*}
\item \emph{Supporting planes give lower bounds only.} Let $z_1=(U_1,p_1,y_1)$ be a point of the cycle, with $p_1'$, $\etaz$, $\zeta_0$, $\kappa_1$ evaluated at $z_1$, and let $z_2=(U_2,p_2,y_2)$ be any point. The quantity \eqref{eq:delta} with base $z_1$ is
\begin{equation}\label{gau:deltakappa}
\Delta=\bigl[\zeta_0(U_2-U_1)-\etaz(p_2-p_1)-(y_2-y_1)\bigr]+\kappa_1\,(p_1,p_1')\wedge(q_2-q_1).
\end{equation}
If $q_2\ne q_1$ lies in the closed region bounded by the projected cycle, then $(p_1,p_1')\wedge(q_2-q_1)<0$, and the supporting inequality $\Delta<0$ is equivalent to the lower bound
\[
\kappa_1>\frac{\zeta_0(U_2-U_1)-\etaz(p_2-p_1)-(y_2-y_1)}{\bigl|(p_1,p_1')\wedge(q_2-q_1)\bigr|}.
\]
This covers every other point of the same cycle, every point of a cycle nested inside it, and $z_2=E$, for which the bound reads $\kappa_1>0$. The inequality $W>0$ reads $\kappa>-W_0/\Xg$.
\end{enumerate}
\end{proposition}
\begin{proof}
(i) The definition of $K$ and the tangent orthogonality \eqref{eq:orthogonality} form the linear system
\[
U\zeta-p\,\eta=K+y-k,\qquad p\,\zeta-p'\eta=x-y
\]
for $(\zeta,\eta)$, with determinant $-Up'+p^2=\Sg\ne0$. Cramer's rule gives
\[
\zeta=\frac{-p'(K+y-k)+p(x-y)}{\Sg}=\frac{\Xg-p'K}{\Sg},\qquad
\eta=\frac{U(x-y)-p(K+y-k)}{\Sg}=\frac{\Sg\etaz-pK}{\Sg},
\]
where $\Xg=p'(k-y)+p(x-y)$ was used. These are the first two formulas in \eqref{gau:affine}; the third is the definition of $\kappa$. Setting $K=0$ in the linear system shows that $(\zeta_0,\etaz)$ solves it with right-hand sides $y-k$ and $x-y$, which are the two stated identities. For $W$, use \eqref{eq:projective} to write $W=p\,[x-(1+\eta)\zeta]-p'\,[\zeta+(\mu-1)\eta-\eta^2]$ and insert the first two formulas. The terms quadratic in $\kappa$ are $-p\cdot pp'\kappa^2+p'\cdot p^2\kappa^2=0$. The coefficient of $\kappa$ is
\[
p\,[(1+\etaz)p'+p\,\zeta_0]+p'\,[p'+(\mu-1)p-2\etaz p]=p'^2+\mu pp'+p\,(\zeta_0p-\etaz p')=p'^2+\mu pp'+p(x-y)=\Xg .
\]

(ii) By \eqref{eq:positiveK}, $K'=k\psiU(U)-(1+\eta)K$, and $\Sg'=(\etaz-\mu)\Sg$ by Theorem~\ref{thm:gauges}(vi). Hence
\[
\kappa'=\frac{K'}{\Sg}-\kappa\,\frac{\Sg'}{\Sg}=\frac{k\psiU(U)}{\Sg}-(1+\eta+\etaz-\mu)\,\kappa ,
\]
and $\eta=\etaz-p\kappa$ gives \eqref{eq:kappariccati}.

(iii) The first row of \eqref{gau:compoundsource} is $\Sg'=-\mu\Sg+\Sg\etaz$ and the third is \eqref{gau:Xistate}. For the second row, $\Sg\etaz=U(x-y)-p(y-k)$ and $y''=xp-(x-y)$ give
\[
(\Sg\etaz)'=U\,[xp-(x-y)]-p'(y-k),
\]
while
\begin{align*}
-\Sg\etaz+\Xg+k\psiU(U)p&=-U(x-y)+p(y-k)-p'(y-k)+p(x-y)+(xU-x+k)p\\
&=-U(x-y)-p'(y-k)+xUp .
\end{align*}
The two agree. Dividing the second and third rows by $\Sg$ and using $\Sg'/\Sg=\etaz-\mu$ gives the defect form.

(iv) By \eqref{eq:exactexponent}, $\lambda_cT=-\int_0^T\kappa\,p\dd$. On the rising arc $U$ increases strictly from $U_{\min}$ to $U_{\max}$ and $p\,\mathrm dt=\mathrm dU$, so that arc contributes $\int_{U_{\min}}^{U_{\max}}\kappa_{\rm up}\,\mathrm dU$. On the falling arc $U$ decreases strictly from $U_{\max}$ to $U_{\min}$, and the contribution is $-\int_{U_{\min}}^{U_{\max}}\kappa_{\rm down}\,\mathrm dU$.

(v) The first identity in \eqref{gau:certificate} is $\eta=\etaz-p\kappa$. The function $\Lambda$ is periodic along the cycle, so $\avg{\Dt\Lambda}=0$. The equivalence is obtained by dividing $\varepsilon_\Lambda-p\kappa\le0$ by $p$ on each arc.

(vi) Insert $\zeta=\zeta_0-p_1'\kappa_1$ and $\eta=\etaz-p_1\kappa_1$ into \eqref{eq:delta}: the terms containing $\kappa_1$ are $\kappa_1[-p_1'(U_2-U_1)+p_1(p_2-p_1)]=\kappa_1(p_1,p_1')\wedge(q_2-q_1)$. The sign of the wedge is \eqref{gau:supportline}. The projection of a cycle nested inside lies in the region bounded by the projected cycle by Theorem~\ref{thm:geometry}, and so does the origin. For $z_2=E$ the bracket in \eqref{gau:deltakappa} equals $-(\zeta_0U_1-\etaz p_1-y_1+k)=0$ by (i), and the wedge equals $-\Sg(z_1)$, so $\Delta=-\kappa_1\Sg=-K$. The last sentence is the fourth formula in \eqref{gau:affine}.
\end{proof}

Part (i) shows that along a cycle the four quantities $(U,p,y,\kappa)$ satisfy a closed autonomous system, and that every function of $(U,p,y,\eta,\zeta,K,W)$ is a function of $(U,p,y,\kappa)$. Parts (v) and (vi) locate the missing estimate. A certificate $\Dt\Lambda+\eta-\mu\le0$ built from a state function requires an upper bound for $\kappa$ on the falling arc and a lower bound on the rising arc. Each supporting-plane inequality of Theorem~\ref{thm:support} with base on the cycle involves $\kappa$ only at its base point and bounds it from below, on both arcs, as long as the second point projects into the region bounded by the cycle. An upper bound for $\kappa_1$ from \eqref{gau:deltakappa} would need a log-bounded entire solution whose projection lies strictly on the outer side of the tangent line at $q_1$; no point of the same cycle, of a cycle nested inside it, or the equilibrium has this property.

\begin{remark}[Upper bounds from the dynamics]\label{gau:dynamicbounds}
Two pointwise upper bounds follow directly from the equations. First, $\zeta<y$ on every cycle: by \eqref{eq:projective} and $y'=x-y$,
\[
(y-\zeta)'+(y-\zeta)=\eta\,\zeta>0,
\]
so the periodic function $y-\zeta$ is positive at its minimum. Second, if $\mu\le\tfrac12$ then $\eta<\ell=\lambda-p$ on every cycle:
\[
(\ell-\eta)'+\mu\,(\ell-\eta)=(y-\zeta)+\eta\,(\eta+1-2\mu)>0 ,
\]
and the same argument applies. With \eqref{gau:affine} and $\zeta>0$ these give $\kappa<\zeta_0/p'$ where $p'>0$, $\kappa<(y-\zeta_0)/|p'|$ where $p'<0$, and, for $\mu\le\tfrac12$, $\kappa<(\ell-\etaz)/|p|$ on the falling arc. We have not been able to derive the sign of $\lambda_c$ from these bounds, and they are not used below.
\end{remark}

\subsection{Polar form and nested cycles}
\begin{proposition}[Polar form]\label{prop:polar}
Put $(U,p)=r(\cos\theta,\sin\theta)$ with $r>0$ and $s=(y-k)/r$. Then $p'=-r(\mu\sin\theta+s)$, and on $\{r>0\}$ system \eqref{eq:logsystem} is
\begin{equation}\label{gau:polarsystem}
\begin{aligned}
(\log r)'&=\sin\theta\,(\cos\theta-\mu\sin\theta-s),\\
\theta'&=-\sigma,\qquad \sigma=\sin^2\theta+\mu\sin\theta\cos\theta+s\cos\theta=\frac{\Sg}{r^2},\\
s'&=\frac{k\,(e^{r\cos\theta}-1)}{r}-s-s\,\sin\theta\,(\cos\theta-\mu\sin\theta-s).
\end{aligned}
\end{equation}
On every cycle $\sigma>0$, so $\theta'<0$, and the cycle is a graph $r=R(\theta)$, $s=s(\theta)$ over $\theta\in\R/2\pi\mathbb Z$ with
\begin{equation}\label{gau:polarorbit}
\frac{\mathrm d\log R}{\mathrm d\theta}=g(s,\theta),\qquad
g(s,\theta)=\frac{\sin\theta\,(s+\mu\sin\theta-\cos\theta)}{\sin^2\theta+\mu\sin\theta\cos\theta+s\cos\theta}.
\end{equation}
The weak exponent is
\begin{equation}\label{gau:polarexponent}
\lambda_cT=-\int_0^{2\pi}\frac{K\,\sin\theta}{r\,\sigma^2}\,\mathrm d\theta
=-\int_0^{2\pi}\frac{K\,r^3\sin\theta}{\Sg^2}\,\mathrm d\theta ,
\end{equation}
where $K$, $r$, $\sigma$, $\Sg$ are evaluated at the point of the cycle with polar angle $\theta$. The fraction $K\sin\theta/(r\sigma^2)$ under the integral sign is positive for $\theta\in(0,\pi)$ and negative for $\theta\in(\pi,2\pi)$.
\end{proposition}
\begin{proof}
From $p'=k-\mu p-y$ we get $p'=-\mu r\sin\theta-rs$. Then
\[
(\log r)'=\frac{UU'+pp'}{r^2}=\frac{Up+pp'}{r^2}=\sin\theta\cos\theta-\sin\theta\,(\mu\sin\theta+s),
\]
\[
\theta'=\frac{Up'-pU'}{r^2}=\frac{Up'-p^2}{r^2}=-\cos\theta\,(\mu\sin\theta+s)-\sin^2\theta=-\frac{\Sg}{r^2}.
\]
Finally $s'=y'/r-s\,(\log r)'$ and $y'=ke^U-y=k(e^{r\cos\theta}-1)-rs$. On a cycle $\Sg>0$ and $\theta$ decreases by exactly $2\pi$ in a period by Theorem~\ref{thm:gauges}(ii), so $t\mapsto\theta$ is a diffeomorphism of $\R/T\mathbb Z$ onto $\R/2\pi\mathbb Z$ that reverses orientation. Dividing the first equation of \eqref{gau:polarsystem} by the second gives \eqref{gau:polarorbit}. For any continuous $T$-periodic $F$, the substitution $\mathrm dt=-\mathrm d\theta/\sigma$ gives $\int_0^TF\dd=\int_0^{2\pi}(F/\sigma)\,\mathrm d\theta$. Apply this to $F=pK/\Sg=K\sin\theta/(r\sigma)$ and use \eqref{eq:exactexponent}. The sign of the integrand is that of $\sin\theta$ because $K,r,\sigma>0$.
\end{proof}

\begin{proposition}[Radial comparison of nested cycles]\label{prop:nestedradii}
Suppose that at some parameter there are two distinct cycles $\Gamma_1$, $\Gamma_2$, with the projection of $\Gamma_1$ inside that of $\Gamma_2$. Let $R_i(\theta)$, $s_i(\theta)$ and $\sigma_i(\theta)=\sin^2\theta+\mu\sin\theta\cos\theta+s_i(\theta)\cos\theta$ be the polar data of $\Gamma_i$ as in Proposition~\ref{prop:polar}. Then $R_2>R_1$ and $s_2>s_1$ for every $\theta$, and
\begin{equation}\label{gau:radialratio}
\frac{\mathrm d}{\mathrm d\theta}\log\frac{R_2}{R_1}=\frac{\sin\theta\,(s_2-s_1)}{\sigma_1\sigma_2}.
\end{equation}
Hence $R_2/R_1$ is strictly increasing on $[0,\pi]$ and strictly decreasing on $[\pi,2\pi]$, and for $\theta\notin\pi\mathbb Z$
\begin{equation}\label{gau:ratiochain}
\frac{U_{\min,2}}{U_{\min,1}}>\frac{R_2(\theta)}{R_1(\theta)}>\frac{U_{\max,2}}{U_{\max,1}}>1 ,
\end{equation}
where $U_{\min,i}$, $U_{\max,i}$ are the trough and crest values of $U$ on $\Gamma_i$. Moreover
\[
\int_0^{2\pi}\frac{\sin\theta\,(s_2-s_1)}{\sigma_1\sigma_2}\,\mathrm d\theta=0 .
\]
The same statements hold for polar coordinates centred at $(-c,0)$ with $0\le c<|U_{\min,1}|$, with $U_{\min,i}+c$ and $U_{\max,i}+c$ in place of $U_{\min,i}$ and $U_{\max,i}$.
\end{proposition}
\begin{proof}
Both projections are strictly star-shaped about the origin (Theorem~\ref{thm:gauges}(ii)), so $R_1,R_2$ are well defined. The point of $\Gamma_1$ at angle $\theta$ lies in the bounded region enclosed by the projection of $\Gamma_2$. The ray at angle $\theta$ meets that projection only at radius $R_2(\theta)$, and its part beyond $R_2(\theta)$ lies in the unbounded component. Hence $R_1(\theta)<R_2(\theta)$.

Fix $\theta$ and let $z_i=(U_i,p_i,y_i)\in\Gamma_i$ be the points at polar angle $\theta$, so $(U_2-U_1,p_2-p_1)=(R_2-R_1)(\cos\theta,\sin\theta)$. Let $\eta,\zeta$ be the adjoint ratios of $\Gamma_1$ at $z_1$ and put $d_1=\zeta\cos\theta-\eta\sin\theta$. Inequality \eqref{eq:positiveK} at $z_1$ is $R_1d_1-(y_1-k)>0$, that is, $d_1>s_1$. Theorem~\ref{thm:support} with base $z_1$ and second solution through $z_2$ gives $(R_2-R_1)d_1-(y_2-y_1)<0$. Therefore
\[
y_2-k=(y_2-y_1)+(y_1-k)>(R_2-R_1)\,d_1+R_1s_1>(R_2-R_1)\,s_1+R_1s_1=R_2s_1 ,
\]
which is $s_2>s_1$.

By \eqref{gau:polarorbit}, $\frac{\mathrm d}{\mathrm d\theta}\log(R_2/R_1)=g(s_2,\theta)-g(s_1,\theta)$. Write $g(s,\theta)=\sin\theta\,(s+n_0)/(s\cos\theta+s_0)$ with $n_0=\mu\sin\theta-\cos\theta$ and $s_0=\sin^2\theta+\mu\sin\theta\cos\theta$. Then
\[
g(s_2,\theta)-g(s_1,\theta)=\frac{\sin\theta\,(s_2-s_1)(s_0-n_0\cos\theta)}{\sigma_1\sigma_2},\qquad s_0-n_0\cos\theta=\sin^2\theta+\cos^2\theta=1 .
\]
This is \eqref{gau:radialratio}. Its right side is positive on $(0,\pi)$ and negative on $(\pi,2\pi)$. The maximum of $R_2/R_1$ is at $\theta=\pi$, where $R_i=|U_{\min,i}|$, and the minimum at $\theta=0$, where $R_i=U_{\max,i}$; the minimum exceeds $1$ because $R_2>R_1$. The integral identity is the integral of the derivative of a $2\pi$-periodic function.

For a centre $(-c,0)$ write $(U+c,p)=r(\cos\theta,\sin\theta)$ and $s=(y-k)/r$. The first two equations of \eqref{gau:polarsystem} hold verbatim with $\sigma=\Sg_c/r^2$, since only $U'=p$ and $p'=k-\mu p-y$ were used. Both cycles are strictly star-shaped about $(-c,0)$ by Theorem~\ref{thm:gauges}(v), because $c<|U_{\min,1}|<|U_{\min,2}|$; the second inequality is the case $\theta=\pi$ of $R_1<R_2$ for the centre at the origin. In the proof of $s_2>s_1$ replace \eqref{eq:positiveK} by $K+c\,\zeta>0$ at $z_1$, which reads $R_1d_1-(y_1-k)>0$ in the shifted coordinates. The rest is unchanged.
\end{proof}
Proposition~\ref{prop:nestedradii} is conditional on the existence of two cycles and derives no contradiction from it. Since $\partial g/\partial s=\sin\theta/\sigma^2$ has the sign of $\sin\theta$, the radial comparison has opposite senses in the two half-planes: by \eqref{gau:radialratio} the ratio $R_2/R_1$ increases with $\theta$ on the upper half-plane and decreases on the lower one. A comparison of a cycle with its radial dilations therefore has no fixed sign. Formally, when the two cycles merge, \eqref{gau:radialratio} linearises to $\mathrm d\delta/\mathrm d\theta=\delta\,K\sin\theta/(r\sigma^2)$ for $\delta=\log(R_2/R_1)$, because $s_2-s_1=\delta\,K/r$ to first order by the two supporting inequalities; the periodicity of $\delta$ then requires the integral in \eqref{gau:polarexponent} to vanish, that is, $\lambda_c=0$. This formal statement is not used.

\subsection{Sign-definite integrals}
\begin{proposition}[Monotone integrals]\label{prop:monotoneintegrals}
On every cycle and for every $(\alpha,\beta)\ne(0,0)$,
\begin{equation}\label{gau:monotone}
\avg{(\alpha p+\beta p')(\alpha\eta+\beta\zeta)}<0 .
\end{equation}
In particular $\avg{\eta p}<0$ and $\avg{\zeta p'}<0$; moreover $\avg{\eta p'}=\avg{\zeta p}$ and $\avg{\eta p}\avg{\zeta p'}>\avg{\zeta p}^2$. For two points of the cycle with equal $U$, the one on the falling arc has the larger $\eta$ and the larger $y$. For two points with equal $p$, the one with $p'<0$ has the larger $U$, the larger $y$ and the larger $\zeta$.
\end{proposition}
\begin{proof}
Fix $(\alpha,\beta)\ne(0,0)$ and put $L(t)=\alpha U(t)+\beta p(t)$ and $g=\alpha\eta+\beta\zeta$. Then $L'=(\alpha,\beta)\cdot(p,p')$ vanishes at exactly two phases in a period (proof of Theorem~\ref{thm:gauges}(iii)). So $L$ has one maximum $L_{\max}$ and one minimum $L_{\min}$ in a period, is strictly increasing from the minimum to the maximum and strictly decreasing afterwards, and each level $m\in(L_{\min},L_{\max})$ is attained at exactly two points: $z_A(m)$, where $L'>0$, and $z_B(m)$, where $L'<0$.

Both points lie on the line $\alpha U+\beta p=m$, so $q_A-q_B=\nu\,(\beta,-\alpha)$ with a real $\nu\ne0$. By \eqref{gau:supportline} at $z_A$ with $q=q_B$,
\[
0>(p_A,p_A')\wedge(q_B-q_A)=-\nu\,\bigl[p_A\cdot(-\alpha)-p_A'\,\beta\bigr]=\nu\,(\alpha p_A+\beta p_A')=\nu\,L'(t_A),
\]
and $L'(t_A)>0$. Hence $\nu<0$.

Theorem~\ref{thm:support} applies to two distinct points of one cycle. With base $z_A$ and with base $z_B$ it gives
\[
\zeta_A(U_B-U_A)-\eta_A(p_B-p_A)-(y_B-y_A)<0,\qquad
\zeta_B(U_A-U_B)-\eta_B(p_A-p_B)-(y_A-y_B)<0 .
\]
Adding,
\begin{equation}\label{gau:gradientmonotone}
(\zeta_A-\zeta_B)(U_A-U_B)-(\eta_A-\eta_B)(p_A-p_B)>0 .
\end{equation}
With $(U_A-U_B,p_A-p_B)=\nu(\beta,-\alpha)$ this is $\nu\,[g_A-g_B]>0$, so $g_A<g_B$. Splitting the period at the extrema of $L$ and using $L$ as the variable on each part,
\[
T\avg{g\,L'}=\int_{L_{\min}}^{L_{\max}}g_A(m)\,\mathrm dm-\int_{L_{\min}}^{L_{\max}}g_B(m)\,\mathrm dm<0 ,
\]
which is \eqref{gau:monotone}. The cases $(1,0)$ and $(0,1)$ are $\avg{\eta p}<0$ and $\avg{\zeta p'}<0$. By \eqref{eq:orthogonality}, $\eta p'=\zeta p-y'$, and $\avg{y'}=0$, so $\avg{\eta p'}=\avg{\zeta p}$. The left side of \eqref{gau:monotone} is then the quadratic form $\alpha^2\avg{\eta p}+2\alpha\beta\avg{\zeta p}+\beta^2\avg{\zeta p'}$, which is negative definite; its determinant is positive.

For $(\alpha,\beta)=(1,0)$ the point $z_A$ is on the rising arc and $z_B$ on the falling arc, $U_A=U_B$, $p_A>p_B$, and $g=\eta$, so $\eta_B>\eta_A$. The two supporting inequalities read $\eta_A(p_A-p_B)<y_B-y_A<\eta_B(p_A-p_B)$, so $y_B>y_A$. For $(\alpha,\beta)=(0,1)$ the point $z_A$ has $p'>0$, $z_B$ has $p'<0$, $p_A=p_B$ and $q_A-q_B=\nu(1,0)$, so $U_B>U_A$; $g=\zeta$ gives $\zeta_B>\zeta_A$, and the two supporting inequalities read $\zeta_A(U_B-U_A)<y_B-y_A<\zeta_B(U_B-U_A)$.
\end{proof}

\begin{corollary}[A convex criterion]\label{gau:convexcriterion}
Let $F(U,p,y,\eta,\zeta)$ be a $C^1$ function on a neighbourhood in $\R^5$ of the curve $t\mapsto(U,p,y,\eta,\zeta)(t)$ of a cycle, and suppose that along this curve
\begin{enumerate}[label=(\arabic*)]
\item for two points with equal $U$, the value of $\partial_UF$ at the point of the rising arc does not exceed its value at the point of the falling arc;
\item for two points with equal $p$, the value of $\partial_pF$ at the point with $p'>0$ does not exceed its value at the point with $p'<0$;
\item $\eta-\mu+\partial_yF\cdot(x-y)+\partial_\eta F\cdot\eta'+\partial_\zeta F\cdot\zeta'\le0$, with strict inequality at some phase.
\end{enumerate}
Then $\lambda_c<0$.
\end{corollary}
\begin{proof}
The function $t\mapsto F$ is periodic, so $\int_0^T\frac{\mathrm d}{\mathrm dt}F\dd=0$, and by the chain rule
\[
\lambda_cT=\int_0^T\Bigl(\eta-\mu+\frac{\mathrm d}{\mathrm dt}F\Bigr)\dd=\oint\partial_UF\,\mathrm dU+\oint\partial_pF\,\mathrm dp+\int_0^T\bigl[\eta-\mu+\partial_yF\,(x-y)+\partial_\eta F\,\eta'+\partial_\zeta F\,\zeta'\bigr]\dd .
\]
As in the proof of Proposition~\ref{prop:monotoneintegrals}, the first integral is the integral over $(U_{\min},U_{\max})$ of the difference in (1) and the second is the integral over the range of $p$ of the difference in (2). Both are nonpositive, and the third is negative by (3).
\end{proof}
The criterion allows the two line integrals to be discarded. Hypothesis (1) holds, by the orderings in Proposition~\ref{prop:monotoneintegrals}, if $\partial_UF$ does not depend on $\zeta$ and is nonincreasing in $p$ and nondecreasing in $y$ and in $\eta$ on a set that contains the coordinate-wise paths between the compared points; similarly for (2) with $\partial_pF$ independent of $\eta$ and nondecreasing in $U$, $y$, $\zeta$. We know no function $F$ that satisfies (1)--(3) on every cycle.

\subsection{Moment identities}
\begin{proposition}[H\"older sandwich]\label{prop:holder}
On every cycle
\begin{equation}\label{gau:holder}
\avg{x^{1/3}}^3<\gamma(\gamma-1)(\gamma-\mu)<k .
\end{equation}
Consequently, with $\rho=k/\kH$ as in \eqref{eq:parameters}, every cycle with $\lambda_c\le0$ satisfies $\avg{e^{U/3}}<\rho^{-1/3}$, and in particular $\avg U<-\log\rho$.
\end{proposition}
\begin{proof}
The functions $u_1=\eta$, $u_2=\zeta/\eta$, $u_3=x/\zeta$ are positive and periodic with $u_1u_2u_3=x$. By \eqref{eq:projective},
\[
(\log\eta)'=\frac{\zeta}{\eta}+\mu-1-\eta,\qquad(\log\zeta)'=\frac{x}{\zeta}-1-\eta .
\]
Averaging, $\avg{u_1}=\gamma-1$, $\avg{u_2}=1+\avg\eta-\mu=\gamma-\mu$ and $\avg{u_3}=1+\avg\eta=\gamma$. Put $\tilde u_i=u_i/\avg{u_i}$. The inequality of arithmetic and geometric means gives pointwise
\[
\frac{x^{1/3}}{\bigl(\avg{u_1}\avg{u_2}\avg{u_3}\bigr)^{1/3}}=(\tilde u_1\tilde u_2\tilde u_3)^{1/3}\le\frac{\tilde u_1+\tilde u_2+\tilde u_3}{3},
\]
with equality only where $\tilde u_1=\tilde u_2=\tilde u_3$. Averaging gives $\avg{x^{1/3}}^3\le\gamma(\gamma-1)(\gamma-\mu)$, and equality would require $\tilde u_1=\tilde u_2=\tilde u_3$ at every phase, by continuity. Then $\zeta=c_1\eta^2$ with a constant $c_1>0$, and the first equation of \eqref{eq:projective} becomes the autonomous scalar equation $\eta'=(c_1-1)\eta^2+(\mu-1)\eta$, whose periodic solutions are constant. A constant $\eta$ makes $\zeta=\eta^2+(1-\mu)\eta$ constant, then $x=(1+\eta)\zeta$ constant, so $p=0$ and the orbit is the equilibrium. Hence the left inequality in \eqref{gau:holder} is strict.

For the right inequality, $\gamma(\gamma-1)(\gamma-\mu)=f_\mu(\avg\eta)$ and $k=\avg x=\avg{f_\mu(\eta)}$ by \eqref{eq:gainprimitive}. Now $f_\mu''(s)=6s+4-2\mu$ is positive for $s>\max(0,\mu-1)$: if $\mu\le2$ this is clear, and if $\mu>2$ then $6s+4-2\mu>4\mu-2>0$. By \eqref{eq:etaminbasic} the values of $\eta$ lie in this interval, on which $f_\mu$ is strictly convex, and $\eta$ is nonconstant. Jensen's inequality is therefore strict: $f_\mu(\avg\eta)<\avg{f_\mu(\eta)}$.

The function $\chi(\gamma)=\gamma(\gamma-1)(\gamma-\mu)$ is strictly increasing for $\gamma>\max(1,\mu)$, and $\gamma=1+\avg\eta$ lies in this range. Since $\chi(A)=\kH$, the condition $\lambda_c\le0$, that is, $\gamma\le A$, gives $\chi(\gamma)\le\kH$. Then \eqref{gau:holder} and $x=ke^U$ give $k\avg{e^{U/3}}^3<\kH$. Finally $e^{\avg U/3}\le\avg{e^{U/3}}$ by Jensen's inequality.
\end{proof}

\begin{proposition}[Lotka--Volterra chain and second-order forms]\label{prop:lvchain}
\begin{enumerate}[label=(\roman*)]
\item On $\Dom$ the functions $r_1=x/y$, $r_2=y/\ell$ and $\ell=\lambda-p$ are positive and satisfy
\begin{equation}\label{gau:lv}
r_1'=r_1\,(\lambda+1-r_1-\ell),\qquad r_2'=r_2\,(\mu-1+r_1-r_2),\qquad \ell'=\ell\,(r_2-\mu).
\end{equation}
The interior equilibrium is $(r_1,r_2,\ell)=(1,\mu,\lambda)$, and on every cycle
\begin{equation}\label{gau:lvmeans}
\avg{\frac xy}=1,\qquad\avg{\frac y\ell}=\mu,\qquad\avg\ell=\lambda .
\end{equation}
\item For every positive $C^3$ function $\varphi$ of $t$ and every constant $c$,
\begin{equation}\label{gau:colehopf}
\mathcal L\Bigl[\frac{\varphi'}{\varphi}+c-1\Bigr]=\frac{(D+c)(D+c-1)(D+c-\mu)\,\varphi}{\varphi},\qquad
\mathcal L[w]=w''+(3w+2-\mu)\,w'+f_\mu(w).
\end{equation}
\item On every cycle the adjoint ratio $\eta$ satisfies
\begin{equation}\label{gau:secondorder}
\eta''+(3\eta+2-\mu)\,\eta'+f_\mu(\eta)=x ,
\end{equation}
and along every solution in $\Dom$, $\mathcal L[y/\ell]=x\,(A+p)/\ell$.
\item The pair $(\tilde\eta,\tilde\zeta)=(y/\ell,\ x/\ell)$ satisfies the first equation of \eqref{eq:projective} exactly and the second up to an explicit defect:
\begin{equation}\label{gau:chainpair}
\tilde\eta'=\tilde\zeta+(\mu-1)\tilde\eta-\tilde\eta^2,\qquad
\tilde\zeta'=x-(1+\tilde\eta)\tilde\zeta+\frac{x\,(A+2p-\lambda)}{\ell},
\end{equation}
and $\avg{\tilde\eta}=\mu$ on every cycle.
\end{enumerate}
\end{proposition}
\begin{proof}
(i) By \eqref{eq:system}, $(\log r_1)'=p-(x-y)/y=(\lambda-\ell)-r_1+1$. Since $k=\mu\lambda$, $\ell'=-p'=y-k+\mu p=y-\mu\ell$, so $(\log\ell)'=r_2-\mu$ and $(\log r_2)'=(x-y)/y-(r_2-\mu)=r_1-1-r_2+\mu$. On a cycle the three logarithms are periodic. Averaging the third equation of \eqref{gau:lv} divided by $\ell$ gives $\avg{r_2}=\mu$, the second then gives $\avg{r_1}=1$, and the first gives $\avg\ell=\lambda$.

(ii) Put $H_3=e^{ct}\varphi$, $H_2=H_3'-H_3$, $H_1=H_2'-\mu H_2$, and $w=H_2/H_3=\varphi'/\varphi+c-1$, $z=H_1/H_3$. By construction $H_3'=H_2+H_3$ and $H_2'=H_1+\mu H_2$, which are the third and second rows of the adjoint system with matrix $B$ of \eqref{eq:adjointmatrices}; the first row, $H_1'=xH_3$, is not assumed. Then
\[
w'=\frac{H_1+\mu H_2}{H_3}-w\,\frac{H_2+H_3}{H_3}=z+(\mu-1)w-w^2,\qquad
z'=\frac{H_1'}{H_3}-z\,(1+w).
\]
Substituting $z=w'+w^2+(1-\mu)w$ into the second relation gives $\mathcal L[w]=H_1'/H_3$: indeed $z'+(1+w)z=w''+2ww'+(1-\mu)w'+(1+w)(w'+w^2+(1-\mu)w)$ and $(1+w)(w^2+(1-\mu)w)=f_\mu(w)$. Finally $H_1'=D(D-\mu)(D-1)H_3=e^{ct}(D+c)(D+c-\mu)(D+c-1)\varphi$.

(iii) For $\eta$, eliminate $\zeta=\eta'+\eta^2+(1-\mu)\eta$ from the second equation of \eqref{eq:projective}; this is the computation in (ii) with $z=\zeta$ and $z'+(1+w)z=x$. For $y/\ell$ take $\varphi=\ell$ and $c=A$ in \eqref{gau:colehopf}: $\varphi'/\varphi+c-1=(y/\ell-\mu)+\mu=y/\ell$, and
\[
(D+A)(D+1)(D+\mu)\,\ell=(D+A)(D+1)\,y=(D+A)\,x=x\,(p+A),
\]
because $(D+\mu)\ell=y$, $(D+1)y=x$ and $x'=px$.

(iv) The first equation in \eqref{gau:chainpair} is the second equation of \eqref{gau:lv}, since $\tilde\zeta=r_1r_2$. For the second, $(\log\tilde\zeta)'=p-(r_2-\mu)$, so
\[
\tilde\zeta'-x+(1+\tilde\eta)\tilde\zeta=\frac x\ell\,\bigl[p-r_2+\mu-\ell+1+r_2\bigr]=\frac{x\,(A+p-\ell)}{\ell}=\frac{x\,(A+2p-\lambda)}{\ell}.
\]
The mean is \eqref{gau:lvmeans}.
\end{proof}
In the chemical variables $r_1=X/Z$, $r_2=\mu Z/Y$ and $\ell=\lambda Y$. System \eqref{gau:lv} is of Lotka--Volterra type with a cyclic chain of interactions of mixed signs; it is not a competitive Lotka--Volterra system. Part (i) is elementary and is not claimed as new. Theorem~\ref{thm:gauges}(vi) and part (iv) exhibit three explicit functions of the state, $\etaz$, $\eta_1$ and $y/\ell$, whose mean over every cycle is exactly $\mu$; the exponent $\lambda_c$ is the mean of the difference between $\eta$ and any of them.

\subsection{Pairing with the strong-stable solution and a comparison bound}
By Lemma~\ref{lem:strongstable}, along every cycle the variational equation $v'=Jv$ has a Floquet solution $e_s=(a,-b,-c)$ with $a,b,c>0$ and multiplier $m_s=e^{-\gamma T}$. In components,
\begin{equation}\label{gau:abc}
a'=-b,\qquad b'=-\mu b-c,\qquad c'=-c-xa ,
\end{equation}
and $\hat a=e^{\gamma t}a$, $\hat b=e^{\gamma t}b$, $\hat c=e^{\gamma t}c$ are positive and $T$-periodic. The ratios $\mathfrak r=b/a$ and $\mathfrak s=c/a$ are positive and $T$-periodic, and \eqref{gau:abc} gives
\begin{equation}\label{gau:rsriccati}
\mathfrak r'=\mathfrak r^2-\mu\,\mathfrak r-\mathfrak s,\qquad \mathfrak s'=\mathfrak r\,\mathfrak s-\mathfrak s-x .
\end{equation}
Nothing else from Section~\ref{sec:weakadjoint} is used here.

\begin{proposition}[Compound pairing and comparison bound]\label{prop:strongidentity}
Let $C$ be the matrix in \eqref{eq:compound}.
\begin{enumerate}[label=(\roman*)]
\item \emph{Pairing.} The vector $\nu=e^{At}(c,b,a)^{\mathsf T}$ solves $\nu'=-C^{\mathsf T}\nu$, and $\hat\nu=e^{\lambda_ct}\nu=(\hat c,\hat b,\hat a)^{\mathsf T}$ is positive and $T$-periodic. For every $T$-periodic $C^1$ vector function $w(t)\in\R^3$,
\begin{equation}\label{gau:pairing}
\lambda_c\,\avg{\hat\nu\cdot w}=-\avg{\hat\nu\cdot(w'-Cw)} .
\end{equation}
\item \emph{Two explicit instances.} With $w=(\ell,y,x)^{\mathsf T}$ and with $w=(1,\mu,\mu)^{\mathsf T}$,
\begin{equation}\label{gau:strongformulas}
\lambda_c=\frac{\avg{\hat a\,x\,(\lambda-A-2p)}}{\avg{\hat c\,\ell+\hat b\,y+\hat a\,x}}
=\frac{\avg{\hat a\,(x-\kH)}}{\avg{\hat c+\mu\hat b+\mu\hat a}} .
\end{equation}
\item \emph{Scalar forms.} $\lambda_c=\avg{\mathfrak r}-A$, and
\begin{equation}\label{gau:chiidentities}
\gamma(\gamma-1)(\gamma-\mu)=\frac{\avg{x\hat a}}{\avg{\hat a}},\qquad k=\avg{\chi(\mathfrak r)},\qquad
\mathfrak r''+(A-3\mathfrak r)\,\mathfrak r'+\chi(\mathfrak r)=x .
\end{equation}
\item \emph{Pointwise bounds.} $\mathfrak r>1$ and $\mathfrak s>(1-\mu)\,\mathfrak r$ at every phase.
\item \emph{Comparison bound.} Let $\varphi,\tilde\eta,\tilde\zeta$ be positive $T$-periodic $C^1$ functions and $w=\varphi\,(1,\tilde\eta,\tilde\zeta)^{\mathsf T}$. Then $w'=Cw+\diag(g_1,g_2,g_3)\,w$ with
\begin{equation}\label{gau:defects}
g_1=(\log\varphi)'+\mu-\tilde\eta,\quad
g_2=(\log\varphi)'+(\log\tilde\eta)'+1-\frac{\tilde\zeta}{\tilde\eta},\quad
g_3=(\log\varphi)'+(\log\tilde\zeta)'+A-\frac{x}{\tilde\zeta},
\end{equation}
and $\lambda_c\le-\avg{\min_ig_i}$. The mean $\avg{\min_ig_i}$ does not depend on $\varphi$: with the residuals $R_1=\tilde\eta'-[\tilde\zeta+(\mu-1)\tilde\eta-\tilde\eta^2]$ and $R_2=\tilde\zeta'-[x-(1+\tilde\eta)\tilde\zeta]$ of \eqref{eq:projective},
\begin{equation}\label{gau:comparison}
\lambda_c\le\avg{\tilde\eta-\mu+\max\Bigl(0,\,-\frac{R_1}{\tilde\eta},\,-\frac{R_2}{\tilde\zeta}\Bigr)},
\end{equation}
with equality for $(\tilde\eta,\tilde\zeta)=(\eta,\zeta)$. If $\log\tilde\eta=\sum_j\vartheta_jF_j$ and $\log\tilde\zeta=\sum_j\varpi_jF_j$ for fixed $C^1$ functions $F_1,\dots,F_m$ of the state, the right side of \eqref{gau:comparison} is a convex function of $(\vartheta,\varpi)\in\R^{2m}$.
\end{enumerate}
\end{proposition}
\begin{proof}
(i) By \eqref{gau:abc},
\[
(e^{At}c)'=e^{At}(\mu c-xa),\qquad(e^{At}b)'=e^{At}(b-c),\qquad(e^{At}a)'=e^{At}(Aa-b),
\]
and the rows of $-C^{\mathsf T}$ are $(\mu,0,-x)$, $(-1,1,0)$, $(0,-1,A)$. Hence $\nu'=-C^{\mathsf T}\nu$. Since $A-\gamma=\mu-\avg\eta=-\lambda_c$, we have $\nu=e^{-\lambda_ct}\hat\nu$. For periodic $w$, $(\nu\cdot w)'=\nu\cdot(w'-Cw)$, that is,
\[
(\hat\nu\cdot w)'-\lambda_c\,\hat\nu\cdot w=\hat\nu\cdot(w'-Cw).
\]
The mean of $(\hat\nu\cdot w)'$ over a period is zero.

(ii) For $w=(\ell,y,x)^{\mathsf T}$: $\ell'-(-\mu\ell+y)=0$ and $y'-(-y+x)=0$ by the equations in the proof of Proposition~\ref{prop:lvchain}(i), and $x'-(x\ell-Ax)=x\,(p-\ell+A)=x\,(2p+A-\lambda)$. So $w'-Cw=(0,0,x(2p+A-\lambda))^{\mathsf T}$ and \eqref{gau:pairing} gives the first formula; the denominator is positive because $\ell>0$ on $\Dom$. For $w=(1,\mu,\mu)^{\mathsf T}$: $Cw=(0,0,x-A\mu)^{\mathsf T}$, so $w'-Cw=(0,0,\kH-x)^{\mathsf T}$, which gives the second.

(iii) From $(\log a)'=-\mathfrak r$ and $a(t+T)=e^{-\gamma T}a(t)$ we get $\avg{\mathfrak r}=\gamma=A+\lambda_c$. By \eqref{gau:abc}, $a''=\mu b+c$ and $a'''=-\mu^2b-Ac-xa$, so $a'''+Aa''+\mu a'+xa=0$. Inserting $a=e^{-\gamma t}\hat a$ gives $[(D-\gamma)^3+A(D-\gamma)^2+\mu(D-\gamma)+x]\hat a=0$. Averaging removes every term that is a derivative of a periodic function and leaves $(-\gamma^3+A\gamma^2-\mu\gamma)\avg{\hat a}+\avg{x\hat a}=0$, and $\gamma^3-A\gamma^2+\mu\gamma=\gamma(\gamma-1)(\gamma-\mu)$. By \eqref{gau:rsriccati},
\[
\bigl(\mathfrak s+\tfrac12\mathfrak r^2-\mathfrak r\bigr)'=\mathfrak r\mathfrak s-\mathfrak s-x+(\mathfrak r-1)(\mathfrak r^2-\mu\mathfrak r-\mathfrak s)=\mathfrak r^3-A\mathfrak r^2+\mu\mathfrak r-x=\chi(\mathfrak r)-x ,
\]
and averaging gives $\avg{\chi(\mathfrak r)}=\avg x=k$. Substituting $\mathfrak s=\mathfrak r^2-\mu\mathfrak r-\mathfrak r'$ into the left side gives $(\mathfrak r-A)\mathfrak r'-\mathfrak r''+2\mathfrak r\mathfrak r'$, which is the third identity in \eqref{gau:chiidentities}.

(iv) Put $d=c-(1-\mu)b$. By \eqref{gau:abc}, $d'=-c-xa+(1-\mu)(\mu b+c)=-\mu d-xa$, so $e^{\mu t}d$ is strictly decreasing. Write $d=e^{-\gamma t}\hat d$ with $\hat d$ periodic; then $e^{(\mu-\gamma)t}\hat d(t)$ is strictly decreasing, and $\gamma>\max(1,\mu)$ by \eqref{eq:etaminbasic}. If $\hat d(t_0)\le0$ for some $t_0$, then
\[
e^{(\mu-\gamma)(t_0+T)}\hat d(t_0+T)=e^{(\mu-\gamma)T}\cdot e^{(\mu-\gamma)t_0}\hat d(t_0)\ge e^{(\mu-\gamma)t_0}\hat d(t_0),
\]
because a nonpositive number is multiplied by a factor in $(0,1)$. This contradicts strict decrease. Hence $d>0$, which is $\mathfrak s>(1-\mu)\mathfrak r$. Next $(b-a)'=-\mu b-c+b=-d<0$, so $b-a=e^{-\gamma t}(\hat b-\hat a)$ is strictly decreasing, and the same argument with $\gamma$ in place of $\gamma-\mu$ gives $b>a$.

(v) The three components of $Cw$ are $\varphi\,(\tilde\eta-\mu)$, $\varphi\,(\tilde\zeta-\tilde\eta)$ and $\varphi\,(x-A\tilde\zeta)$. Subtracting them from $\varphi'$, $(\varphi\tilde\eta)'$, $(\varphi\tilde\zeta)'$ and dividing by $\varphi$, $\varphi\tilde\eta$, $\varphi\tilde\zeta$ gives \eqref{gau:defects}. Direct subtraction gives $g_2-g_1=R_1/\tilde\eta$ and $g_3-g_1=R_2/\tilde\zeta$, hence
\[
\min_ig_i=(\log\varphi)'+\mu-\tilde\eta-\max\Bigl(0,-\frac{R_1}{\tilde\eta},-\frac{R_2}{\tilde\zeta}\Bigr),
\]
and $\avg{(\log\varphi)'}=0$. It remains to prove $\lambda_c\le-\avg m$ for $m=\min_ig_i$, which is continuous and periodic. Put $\mathcal E(t)=\exp(-\int_0^tm)$ and $\bar w=\mathcal Ew$. Then $\bar w'=C\bar w+\diag(g_i-m)\bar w$. The off-diagonal entries of $C$ are nonnegative, so the evolution operator $\Phi_C(t,s)$ of $w'=Cw$ is a nonnegative matrix for $t\ge s$. The variation-of-constants formula gives
\[
\bar w(T)=\Phi_C(T,0)\bar w(0)+\int_0^T\Phi_C(T,s)\,\diag\bigl(g_i(s)-m(s)\bigr)\,\bar w(s)\,\mathrm ds\ \ge\ \Phi_C(T,0)\bar w(0)
\]
componentwise, because $\bar w>0$ and $g_i-m\ge0$. Since $\bar w(0)=w(0)$ and $\bar w(T)=e^{-T\avg m}w(0)$,
\[
\Phi_C(T,0)\,w(0)\le e^{-T\avg m}\,w(0),\qquad w(0)>0 .
\]
A nonnegative matrix $P$ has a nonnegative left eigenvector $u\ne0$ for the eigenvalue $\spr P$; pairing it with $Pw(0)\le e^{-T\avg m}w(0)$ and using $u\cdot w(0)>0$ gives $\spr\Phi_C(T,0)\le e^{-T\avg m}$.

It remains to identify $\spr\Phi_C(T,0)$. For two solutions $v^1,v^2$ of $v'=Jv$, direct differentiation shows that
\[
\bigl(v^1_1v^2_2-v^1_2v^2_1,\ \ v^1_3v^2_1-v^1_1v^2_3,\ \ v^1_3v^2_2-v^1_2v^2_3\bigr)^{\mathsf T}
\]
solves $w'=Cw$; this is the sign-conjugated wedge used for \eqref{eq:compound}. The map $(v^1,v^2)\mapsto w$ is bilinear, alternating and onto $\R^3$, so $\Phi_C(T,0)$ represents the operator induced on the exterior square of $\R^3$ by the monodromy matrix of $v'=Jv$. Its eigenvalues are the pairwise products of the multipliers $1,m_s,m_t$, counted with multiplicity: $m_s$, $m_t$, $m_sm_t$. By \eqref{eq:multipliers} the largest is $m_t=e^{\lambda_cT}$. Therefore $e^{\lambda_cT}\le e^{-T\avg m}$.

For the true pair, $R_1=R_2=0$ and the right side of \eqref{gau:comparison} is $\avg\eta-\mu=\lambda_c$. For the last assertion take $\varphi=1$. With $\mathbb F=(F_1,\dots,F_m)$,
\[
g_1=\mu-e^{\vartheta\cdot\mathbb F},\qquad
g_2=\vartheta\cdot\Dt\mathbb F+1-e^{(\varpi-\vartheta)\cdot\mathbb F},\qquad
g_3=\varpi\cdot\Dt\mathbb F+A-x\,e^{-\varpi\cdot\mathbb F}
\]
at each phase. Each is an affine function of $(\vartheta,\varpi)$ minus a positive multiple of the exponential of a linear function, hence concave. A pointwise minimum of concave functions is concave, and so is its mean; the right side of \eqref{gau:comparison} is minus that mean.
\end{proof}
The second formula in \eqref{gau:strongformulas} and the first identity in \eqref{gau:chiidentities} express the same fact: $\lambda_c<0$ if and only if the $\hat a$-weighted mean of $x$ is smaller than $\kH$. For $w=\varphi(1,\tilde\eta,\tilde\zeta)^{\mathsf T}$ the pairing \eqref{gau:pairing} reads $\lambda_c\avg{\hat\nu\cdot w}=-\avg{\sum_ig_i\hat\nu_iw_i}$, an exact weighted form of the comparison bound. For the explicit pair of Proposition~\ref{prop:lvchain}(iv) one has $R_1=0$ and $-R_2/\tilde\zeta=\lambda-A-2p$, so \eqref{gau:comparison} gives only $\lambda_c\le\avg{\max(0,\lambda-A-2p)}$, which is positive in $\cP$. Jensen's inequality applied to $k=\avg{\chi(\mathfrak r)}$ gives $\gamma<\gamma_E$, which is weaker than $\gamma<A$.

\subsection{The Gauss map}
\begin{proposition}[Injectivity of the Gauss map]\label{prop:gaussmap}
Fix a parameter and let $z_1\ne z_2$ be points of cycles, of the same cycle or of two different cycles, or let one of them be $E$. Then their adjoint ratios differ: $(\eta,\zeta)(z_1)\ne(\eta,\zeta)(z_2)$, where $(\eta,\zeta)(z)$ denotes the adjoint ratios at the point $z$ of its cycle and the ratios of $E$ are $(\eta_E,\zeta_E)$. Consequently the image of each cycle under $t\mapsto(\eta,\zeta)$ is a regular simple closed curve, the images of distinct cycles are disjoint, and none passes through $(\eta_E,\zeta_E)$. In the coordinates $(\eta,\omega)$, $\omega=\zeta+\eta+\eta^2/2$, in which
\begin{equation}\label{gau:lienardplane}
\eta'=\omega-G_\mu(\eta),\qquad \omega'=x-f_\mu(\eta),\qquad G_\mu(s)=\tfrac32s^2+(2-\mu)s,
\end{equation}
each of these curves is traversed clockwise.
\end{proposition}
\begin{proof}
Let $n_i=\bigl(\zeta(z_i),-\eta(z_i),-1\bigr)$ for $i=1,2$. Theorem~\ref{thm:support} with base $z_1$ gives $n_1\cdot(z_2-z_1)<0$, and with base $z_2$ it gives $n_2\cdot(z_1-z_2)<0$. If $n_1=n_2$, adding the two inequalities gives $0<0$. Hence the map from a cycle to the $(\eta,\zeta)$ plane is injective, and images of distinct cycles and of $E$ are disjoint. The curve is regular because $W=\zeta'p-\eta'p'>0$ excludes $\eta'=\zeta'=0$.

The first equation in \eqref{gau:lienardplane} is the first equation of \eqref{eq:projective} with $\zeta=\omega-\eta-\eta^2/2$, and the second is \eqref{eq:gainprimitive}. The map $(\eta,\zeta)\mapsto(\eta,\omega)$ is a diffeomorphism of the plane with Jacobian determinant $1$, so the image in the $(\eta,\omega)$ plane is again a regular simple closed curve. Let $(\eta_*,\omega_{\max})$ be a point of the curve with maximal $\omega$. There $\omega'=0$, so $\eta'\ne0$ by regularity. The function $G_\mu$ is strictly increasing on $s>\max(0,\mu-1)$, since $G_\mu'(s)=3s+2-\mu$ is positive there, and the values of $\eta$ lie in that interval by \eqref{eq:etaminbasic}. At the phase where $\eta$ is maximal, $\eta'=0$, so $\omega=G_\mu(\eta_{\max})\ge G_\mu(\eta_*)$, and $\omega\le\omega_{\max}$. Hence $\omega_{\max}\ge G_\mu(\eta_*)$, that is, $\eta'\ge0$ at the top point, and therefore $\eta'>0$ there. At a point of maximal ordinate the curve lies in the closed half-plane below its horizontal tangent, so the bounded complementary component of the Jordan curve lies locally below that tangent. The velocity there points in the direction of increasing abscissa, so the interior is on the right of the velocity. A regular simple closed curve with its interior on the right of its velocity at one point is traversed clockwise.
\end{proof}

\subsection{What the identities do not decide}
None of the identities of this section determines the sign of $\lambda_c$. Every exact formula for the exponent in Theorem~\ref{thm:gauges} and Propositions~\ref{prop:kappa} and~\ref{prop:polar} has the form $\lambda_c=-\avg{pN/S}$ with $N>0$ and $S>0$, and the factor $p$ is positive on the rising arc and negative on the falling arc. The gauges therefore bound $\eta-\mu$ from above on the rising arc and from below on the falling arc, and the supporting planes bound $\kappa$ from below on both; a proof of $\lambda_c<0$ by these means would need an upper bound for $\kappa$, equivalently for $\eta$ or $K$, on the falling arc. Inequality \eqref{gau:holder}, equation \eqref{gau:secondorder}, the second formula in \eqref{gau:strongformulas} and the identities \eqref{gau:chiidentities} use only equations \eqref{eq:projective} and \eqref{gau:abc}, in which $x$ enters as a coefficient; they remain valid, with $\le$ in \eqref{gau:holder}, when $x$ is replaced by any positive periodic function of mean $k$, including the constant $k$, for which $\avg\eta-\mu=\gamma_E-A>0$ in $\cP$. These identities are thus compatible with either sign.
